\documentclass[11pt]{article}
\usepackage[a4paper,margin=27mm]{geometry}
\usepackage[T1]{fontenc}
\usepackage[utf8]{inputenc}
\usepackage{lmodern}
\usepackage{amsmath,amssymb,amsthm,bm}
\usepackage{booktabs}
\usepackage{microtype}
\usepackage[colorlinks=true,linkcolor=blue,citecolor=blue,urlcolor=blue]{hyperref}

\newtheorem{theorem}{Theorem}[section]
\newtheorem{remark}[theorem]{Remark}
\newcommand{\F}{\mathbb F_2}
\newcommand{\dist}{\operatorname{dist}}

\title{Understanding the Phase-1 Theory:\\
A Swim-Distance Construction for the Concatenated Color-Code Decoder}
\author{Explanatory research note}
\date{September 2026}

\begin{document}
\maketitle

\begin{abstract}
This note gives a streamlined explanation of the Phase-1 theory developed for
a swim-distance-like soft output for the concatenated matching decoder of a
triangular color code. The aim is not to reproduce every combinatorial lemma,
but to make the logical structure transparent to a reader familiar with the
concatenated matching decoder of Lee, Li, and Bartlett and the surface-code
soft-output construction of Meister, Pattison, and Preskill.

The central observation is that the ordinary stage-2 monochromatic graph
$L_c^*$ merges two geometrically different terminations into one artificial
matching boundary vertex. Once these two terminations are resolved, the
resulting graph has two virtual terminals. For the standard odd-distance
triangular 6.6.6 color-code family considered here, a chain joining these two
terminals maps to a nontrivial physical logical operator, whereas a closed
difference chain within the same stage-2 fiber is stabilizer-trivial. This
gives the logical path structure required to transplant the Meister-type
cluster-contraction construction to each fixed color branch.

The precise Phase-1 scope is important: the construction is for one fixed
color branch, one CSS sector, perfect syndrome measurements, and a fixed
stage-1 output. It is therefore not yet a soft output for the complete
three-branch concatenated decoder, nor is it a syndrome-conditioned
logical-class likelihood ratio.
\end{abstract}

\section{What Phase 1 establishes}

For the planar surface code, a nontrivial logical operator can be represented
by a path joining two inequivalent boundaries. Meister, Pattison, and Preskill
use this geometry to define a decoder confidence: grow the decoder clusters,
contract the regions already covered by those clusters, and measure the
remaining shortest distance between the two logical boundaries \cite{Meister}.

For a triangular color code, the physical code is not itself a simple
two-boundary path problem. A single-qubit Pauli error can excite checks of
three colors, and logical operators admit colored-string and string-net
representatives. Moreover, the concatenated matching decoder does not decode
the full color-code hypergraph in one matching problem. It first projects
information to a restricted graph and then solves a second, monochromatic
matching problem \cite{Lee}.

The Phase-1 question is therefore narrower:
\begin{quote}
For a fixed color $c$, does the \emph{stage-2} monochromatic graph contain a
two-boundary logical topology analogous to the planar surface code? If so,
can Meister's cluster-distance construction be applied to that graph?
\end{quote}

The answer established in Phase 1 is yes, after one essential modification:
the single artificial boundary vertex of the ordinary stage-2 graph must be
resolved into two distinct terminal types.

\section{Color-code ingredients}

\subsection{Triangular color code}

Consider the standard triangular 6.6.6 color-code patch of odd distance
$d\ge 3$. Data qubits live on vertices of a trivalent, three-face-colorable
lattice. Faces have colors $c\in\{r,g,b\}$. For each face $f$, the CSS
stabilizers are
\[
  S_f^X=\prod_{q\in f}X_q,
  \qquad
  S_f^Z=\prod_{q\in f}Z_q.
\]
The triangular patch has three physical boundary sides
$\Gamma_r,\Gamma_g,\Gamma_b$. A boundary of color $c$ is adjacent only to
faces of the other two colors. The code encodes one logical qubit, and a
complete physical side supports representatives of both logical $X$ and
logical $Z$ \cite{Bombin,Lee}.

For the discussion below we analyze $X$ errors using measured $Z$ checks. The
$Z$-error case follows by exchanging the two CSS sectors.

\subsection{Colored strings versus the stage-2 graph}

A physical $c$-colored Pauli string can be built from elementary two-qubit
operators acting on endpoints of primal edges of color $c$. Interior
syndromes cancel along the string, leaving excitations only at its ends.
Color-code strings can also meet to form string nets because of the fusion
relations among the three color charges \cite{Kesselring}.

This physical string picture is useful background, but it should not be
identified directly with the concatenated decoder's stage-2 graph. The
Phase-1 result is instead a statement about how chains in that stage-2 graph
map back to physical-qubit supports.

\section{The concatenated matching decoder}

Fix one color $c$ and denote the other two colors by $c_1,c_2$.

\subsection{Stage 1: restricted graph}

The first graph is $L_{\neg c}^*$. Its real vertices are checks of colors
$c_1$ and $c_2$, together with an artificial matching boundary vertex. Each
primal color-$c$ edge corresponds bijectively to one stage-1 graph edge,
\[
  \epsilon_{\neg c}:E_c\longrightarrow\Delta_1(L_{\neg c}^*).
\]
Matching the measured non-$c$ syndrome gives a set of predicted primal
$c$-edge parities, denoted $E_{\mathrm{pred}}^{(c)}$.

The output of stage 1 is not yet a physical-qubit correction. It becomes part
of the vertex syndrome used by stage 2.

\subsection{Stage 2: monochromatic graph}

The second graph is $L_c^*$. Its real vertices have two different origins,
\[
  \Delta_0(L_c^*)_{\rm real}=F_c\sqcup E_c.
\]
Thus a stage-2 real vertex is either a physical color-$c$ face/check or a
primal color-$c$ edge carrying the auxiliary parity predicted by stage 1.
There is also one artificial matching boundary vertex $v_{\mathrm{bdry}}$.

The essential bijection is
\[
  \epsilon_c:Q\longrightarrow\Delta_1(L_c^*),
\]
where $Q$ is the set of physical data qubits. Every stage-2 graph edge
therefore represents exactly one physical data qubit. For a bulk qubit, the
edge joins the unique incident $c$-face to the unique incident primal
$c$-edge. Near a boundary, one of these two objects may be absent, in which
case the graph edge is attached to $v_{\mathrm{bdry}}$ \cite{Lee}.

Stage 2 solves the matching problem with input
\[
  \sigma_Z^{(c)}\sqcup E_{\mathrm{pred}}^{(c)},
\]
and $\epsilon_c^{-1}$ maps the selected stage-2 graph edges to a physical
qubit correction candidate $V_{\mathrm{pred}}^{(c)}$.

The ordinary decoder repeats both stages for $c=r,g,b$ and chooses the
minimum-weight candidate among the three branches \cite{Lee}.

\section{The hidden boundary structure of $L_c^*$}

For a fixed color $c$, the graph edges incident on $v_{\mathrm{bdry}}$ do not
all represent the same geometric termination. For the standard triangular
family, they split exactly into two classes.

\paragraph{Class 0: the complete $c$-colored side.}
Every qubit on the physical side $\Gamma_c$, including its two corners, has
no incident physical $c$-face. Its stage-2 edge therefore has one real
endpoint, a primal $c$-edge vertex, and one missing face endpoint.

\paragraph{Class 1: the corner opposite $\Gamma_c$.}
The unique corner at which the other two physical sides meet has an incident
$c$-face but no incident primal $c$-edge. Its stage-2 edge therefore has one
real endpoint, a $c$-face vertex, and one missing edge endpoint.

The ordinary graph connects both classes to the same $v_{\mathrm{bdry}}$.
This is valid for matching because the virtual boundary has no prescribed
parity. However, the identification hides the distinction needed for logical
topology.

\subsection{Boundary-resolved graph}

Define the resolved graph $\widetilde L_c^*$ by replacing
$v_{\mathrm{bdry}}$ with two virtual terminals
\[
  b_c^0,\qquad b_c^1.
\]
The $d$ qubit edges associated with the complete side $\Gamma_c$ terminate at
$b_c^0$, while the unique edge associated with the opposite corner terminates
at $b_c^1$. All other edges and real vertices are unchanged. Hence
\[
  \Delta_0(\widetilde L_c^*)
   =F_c\sqcup E_c\sqcup\{b_c^0,b_c^1\},
\]
and there remains a bijection
\[
  \widetilde\epsilon_c:
  Q\longrightarrow\Delta_1(\widetilde L_c^*).
\]

Identifying $b_c^0$ and $b_c^1$ again recovers the ordinary graph $L_c^*$.
Therefore the resolution does not change the physical edge set or its edge
weights. Operationally, ordinary decoding may remain on $L_c^*$; the resolved
graph is used for logical-topology and soft-output analysis.

\section{From stage-2 chains to physical Pauli operators}

Because stage-2 edges are in bijection with physical qubits, any binary graph
chain $\Gamma\subseteq\Delta_1(\widetilde L_c^*)$ defines a physical support
\[
  T_c(\Gamma)=\widetilde\epsilon_c^{-1}(\Gamma).
\]
For the $X$ sector, the associated Pauli operator is
\[
  P_{X,c}(\Gamma)=\prod_{q\in T_c(\Gamma)}X_q.
\]
Addition of graph chains is symmetric difference. Since $X_q^2=I$,
\[
  P_{X,c}(\Gamma_1+\Gamma_2)
  =P_{X,c}(\Gamma_1)P_{X,c}(\Gamma_2).
\]
Thus a difference between two stage-2 corrections is naturally represented
by a binary chain in the resolved graph.

\subsection{What the graph boundary means}

The real stage-2 graph boundary contains two kinds of coordinates:
\begin{enumerate}
  \item parity at a $c$-face vertex is the physical syndrome on that
        $c$-colored check;
  \item parity at a primal $c$-edge vertex is the auxiliary stage-1 parity.
\end{enumerate}
The second quantity is not itself a physical check. However, the collection
of $c$-edge parities determines the syndrome on the non-$c$ physical faces by
the same restricted incidence used in stage 1. Consequently the physical
syndrome factors through the real stage-2 incidence,
\[
  H\,T_c(\Gamma)=\Lambda_c D_c\Gamma,
\]
for an appropriate reconstruction map $\Lambda_c$. The virtual terminals are
not measured checks and are removed from $D_c$.

Therefore a chain whose graph boundary is supported only on the two virtual
terminals has zero physical syndrome. In particular,
\[
  \partial\Gamma=b_c^0+b_c^1
  \quad\Longrightarrow\quad
  H\,T_c(\Gamma)=0.
\]
This makes a terminal-to-terminal path a candidate logical or stabilizer
operator. Zero physical syndrome alone does not yet decide which.

\section{Why terminal-to-terminal means logical}

This is the central topological statement of Phase 1.

Consider differences between corrections with the same real stage-2 input.
Such a difference $z$ has no boundary on any real stage-2 vertex, so its graph
boundary can only be
\[
  \partial z=0
  \qquad\text{or}\qquad
  \partial z=b_c^0+b_c^1.
\]

To determine its physical logical class, use a known physical logical $Z$
representative supported on the complete side $\Gamma_c$. The key parity
identity proved in Phase 1 is
\[
  \text{overlap of }T_c(z)\text{ with }\Gamma_c
  =
  (\partial z)_{b_c^0}
  \pmod 2.
\]
This holds because exactly the qubit edges belonging to $\Gamma_c$ terminate
at $b_c^0$.

The triangular code encodes one qubit. Hence a zero-syndrome $X$-type
operator is stabilizer-trivial exactly when it commutes with a nontrivial
logical $Z$, and it is in the nontrivial logical-$X$ class when it
anticommutes with that logical $Z$. Therefore,
\[
  \boxed{
  \partial z=0
  \Longrightarrow
  T_c(z)\text{ is stabilizer-trivial},
  }
\]
while
\[
  \boxed{
  \partial z=b_c^0+b_c^1
  \Longrightarrow
  T_c(z)\text{ is a nontrivial logical }X.
  }
\]

The full Phase-1 proof additionally shows that every closed resolved chain in
the fixed stage-2 fiber is generated by the appropriate physical face
stabilizers, and that every nontrivial fixed-fiber logical class has a simple
$b_c^0$--$b_c^1$ path representative. Those details are appropriate for a
formal appendix; the parity argument above is the main conceptual reason.

\begin{theorem}[Fixed-branch logical topology, streamlined form]
For the standard odd-distance triangular 6.6.6 family, fixed color $c$, and
one CSS sector, the relative difference space of a fixed stage-2 input has
exactly two physical equivalence classes. Closed chains represent the trivial
class, while chains joining $b_c^0$ to $b_c^1$ represent the nontrivial
logical class.
\end{theorem}

\begin{remark}
The phrase ``fixed stage-2 input'' is essential. The auxiliary parity on the
$E_c$ vertices comes from stage 1. Phase 1 classifies logical differences
\emph{within that fixed auxiliary-parity fiber}. It does not say that stage 1
is correct, nor does it compare different stage-1 predictions.
\end{remark}

\section{Why the Meister construction now applies}

After the theorem above, the relevant logical structure of
$\widetilde L_c^*$ is the same structure used by Meister for a planar
surface-code soft output:
\begin{itemize}
  \item graph edges carry physical error labels and nonnegative weights;
  \item two inequivalent virtual boundaries remain distinct;
  \item a chain between the two boundaries is nontrivial logical;
  \item a same-syndrome closed difference is stabilizer-trivial.
\end{itemize}

This does \emph{not} mean that the color code is a surface code, or that the
complete concatenated decoder is equivalent to a surface-code decoder. In
particular, the real vertices of $\widetilde L_c^*$ have mixed meaning:
$F_c$ vertices are physical checks whereas $E_c$ vertices are stage-1
auxiliary constraints.

The statement is only that the fixed-branch stage-2 logical quotient has the
two-terminal path topology required for the Meister construction.

\section{Decoder clusters and contraction}

In Meister's MWPM construction, clusters are obtained from the dual growth of
the matching problem, not from connected components of the final correction
support \cite{Meister}. Conceptually, syndrome vertices grow metric balls with
decoder-determined radii. The union of these balls forms a covered region
$\mathcal C_c$, and its connected components are the decoder clusters.

For the present construction, this growth is interpreted on the resolved
graph $\widetilde L_c^*$. The two terminals remain distinct; touching one
terminal must not create a virtual zero-cost shortcut through the other.

Each connected component of $\mathcal C_c$ is contracted to zero metric
diameter. Equivalently, distance is charged only on portions of graph edges
not covered by the grown clusters. If an edge $e=uv$ has original weight
$w_e$ and cluster coverage extends distances $h_u$ and $h_v$ into the edge
from its two endpoints, then the residual weight is
\[
  \bar w_e=\max\{0,w_e-h_u-h_v\}.
\]
This partial-edge treatment matters: assigning a whole edge zero cost merely
because both endpoints belong to one cluster can be incorrect.

\section{The per-color swim distance}

Define
\[
  \boxed{
  \phi_c=\dist_{\bar G_c}(b_c^0,b_c^1),
  }
\]
where $\bar G_c$ is the resolved stage-2 graph with cluster-contracted
residual edge weights.

Because terminal-to-terminal paths are exactly the nontrivial fixed-fiber
logical class,
\[
  \phi_c
  =
  \min_{\substack{
       z\ \text{same-fiber difference}\\
       z\ \text{nontrivial logical}}}
       \bar w_c(z).
\]
Equivalently, it is the minimum residual cost of a simple
$b_c^0$--$b_c^1$ path.

Thus the surface-code intuition survives: the decoder is more confident when
a larger amount of uncovered metric length is required to complete a
nontrivial logical path. The nontrivial color-code step is the proof that the
two resolved stage-2 terminal types really distinguish physical logical
parity.

\section{A Meister-type representative-gap bound}

For exact stage-2 MWPM with the certified optimal dual used to define the
clusters, Phase 1 proves a one-sided analogue of Meister's
minimum-representative bound. Let $f_c$ be the stage-2 MWPM correction in the
fixed stage-1 fiber, and define
\[
  W_{c,\mathrm{base}}^{(2)}=w_c(f_c).
\]
Let $W_{c,\mathrm{opp}}^{(2)}$ be the minimum weight of a correction with the
same real stage-2 input but in the opposite physical logical class. Then
\[
  \boxed{
  W_{c,\mathrm{opp}}^{(2)}-W_{c,\mathrm{base}}^{(2)}\ge\phi_c.
  }
\]

The logic is the same as in Meister's proof: the optimal dual certifies that
the MWPM correction is fully covered by the grown clusters, while any
opposite-class correction differs from it by a nontrivial terminal-joining
chain whose uncovered cost is at least $\phi_c$.

This theorem requires the appropriate optimal matching dual. Arbitrary radii,
a suboptimal dual, or generic cluster information do not imply the bound.

\section{What $\phi_c$ does and does not mean}

Three statements should be kept separate.

\paragraph{Exact geometric statement.}
$\phi_c$ is exactly the minimum cluster-contracted cost of a nontrivial
logical difference in the fixed color-$c$, fixed-stage-1 stage-2 problem.

\paragraph{Certified minimum-representative statement.}
For exact MWPM with the specified optimal dual,
\[
  W_{c,\mathrm{opp}}^{(2)}-W_{c,\mathrm{base}}^{(2)}\ge\phi_c.
\]

\paragraph{Probability statement.}
No exact logical-class likelihood ratio is established. Even with independent
physical bit-flip weights, a logical-class posterior contains a sum over many
degenerate representatives. The complete concatenated decoder also contains
other stage-1 fibers and two other color branches. Therefore Phase 1 does not
identify $\phi_c$ with the full syndrome-conditioned logical LLR, with the
comparative gap of the complete decoder, or with the final logical failure
probability.

\section{Computation}

Assume the stage-2 graph, physical edge weights, and decoder growth radii are
already available. A direct computation is:
\begin{enumerate}
  \item restore the two resolved terminals $b_c^0,b_c^1$ from the labelled
        ordinary stage-2 graph;
  \item use the decoder radii to determine the amount of cluster coverage
        reaching every graph vertex;
  \item compute $\bar w_e=\max\{0,w_e-h_u-h_v\}$ for every edge;
  \item run a shortest-path search from $b_c^0$ to $b_c^1$ in the residual
        graph.
\end{enumerate}

The coverage values $h_v$ can be computed by one additional multi-source
shortest-path pass. Thus, with the radii already supplied, the postprocessing
requires two Dijkstra-type passes and linear edge work,
\[
  O\!\left((|V|+|E|)\log |V|\right),
\]
with $O(|V|+|E|)$ storage. Since the triangular stage-2 graph has
$|V|,|E|=O(d^2)$, the per-color postprocessing is
\[
  O(d^2\log d)
\]
with standard sparse-graph data structures. Running all three colors changes
only the constant factor.

\subsection{What differs from the surface-code case}

The asymptotic shortest-path complexity is essentially unchanged. The main
differences are structural and implementation-related:
\begin{enumerate}
  \item the ordinary color-code stage-2 graph first requires the
        side-versus-opposite-corner boundary resolution;
  \item part of the stage-2 vertex data are auxiliary stage-1 parity
        constraints rather than directly measured checks;
  \item the resulting confidence is conditional on one fixed stage-1 fiber;
  \item an exact MWPM representative-gap certificate needs matching
        dual/growth information, not only the final correction.
\end{enumerate}

In the inspected current implementation, ordinary PyMatching-style decoding
returns correction and weight information but does not directly expose the
optimal dual radii required by the certified construction. Decoder
instrumentation or a separate certificate-producing procedure is therefore a
separate implementation task.

\section{The Phase-1 logic in one sequence}

The construction can be summarized as
\[
\boxed{
\begin{array}{c}
\text{triangular color code}\\[1mm]
\downarrow\\
\text{concatenated decoder, fixed color }c\\[1mm]
\downarrow\\
L_c^*\text{ with one merged matching boundary}\\[1mm]
\downarrow\ \text{resolve physical provenance}\\
\widetilde L_c^*\text{ with }b_c^0,b_c^1\\[1mm]
\downarrow\ \text{edge}\leftrightarrow\text{qubit map}\\
b_c^0\text{--}b_c^1\text{ chain}
\leftrightarrow
\text{nontrivial physical logical}\\[1mm]
\downarrow\ \text{decoder cluster contraction}\\
\phi_c=\dist_{\bar G_c}(b_c^0,b_c^1).
\end{array}}
\]

The difficult mathematical part is concentrated in the middle: proving that
resolved terminal parity equals physical logical parity. Once that
correspondence is established, the swim-distance construction itself becomes
the same graph-metric idea used for the surface code.

\section{What remains beyond Phase 1}

Phase 1 deliberately does not answer:
\begin{itemize}
  \item how to combine $\phi_r,\phi_g,\phi_b$ after the ordinary decoder
        selects one of the three branches;
  \item how stage-1 uncertainty or failure should enter the confidence;
  \item how the result relates to a full-decoder forced, comparative, or
        exact logical gap;
  \item how well $\phi_c$ calibrates logical failure probability;
  \item how the construction should be extended to circuit-level spacetime
        decoding;
  \item how the principle generalizes beyond matching, for example to a
        direct hypergraph or BP--LSD decoder.
\end{itemize}

Keeping these questions separate is important. Phase 1 establishes the
stage-2 logical geometry and a valid Meister-type soft-output construction,
not the final confidence measure of the complete color-code decoder.

\begin{thebibliography}{9}

\bibitem{Lee}
S.-H. Lee, A. Li, and S. D. Bartlett,
\emph{Color code decoder with improved scaling for correcting circuit-level noise},
Quantum \textbf{9}, 1609 (2025).

\bibitem{Meister}
N. Meister, C. A. Pattison, and J. Preskill,
\emph{Efficient soft-output decoders for the surface code},
arXiv:2405.07433.

\bibitem{Bombin}
H. Bomb\'in and M. A. Mart\'in-Delgado,
\emph{Topological Quantum Distillation},
Phys. Rev. Lett. \textbf{97}, 180501 (2006).

\bibitem{Kesselring}
M. S. Kesselring, F. Pastawski, J. Eisert, and B. J. Brown,
\emph{The boundaries and twist defects of the color code and their applications to topological quantum computation},
Quantum \textbf{2}, 101 (2018).

\end{thebibliography}

\end{document}
